%% 4-4 三平方の定理（第13回：新しく書いた。組 V4-3。ページ制限なし。偶数ページ）
%% ほかの巻が指す形：
%%   3-2「A(−2, 1) と C(3, −2) のように斜めに並んだ2点の距離は、4-4 で」→ はてな？とその答え（√34）
%%   2-8「正方形の対角線が √2」→ ④で回収。2-9「12角形から先の数は、4-4 の三平方の定理をくり返し使うと計算できる」→ もっと知りたい
%%   5-16「三平方の定理（4-4）から x²＋y²＝r² より y＝√(r²−x²)」→ ⑦
\makeatletter
% 本文の中の小見出し（それだけで1行）  \bkivehead{見出し}
\newcommand\bkivehead[1]{\par\addvspace{4pt}\Needspace{4\baselineskip}\noindent{\sffamily\bfseries #1}\par\nobreak\noindent\ignorespaces}
% 式を、まん中に1行で置く  \bkiveline{…}
\newcommand\bkiveline[1]{\par\addvspace{3pt}{\centering #1\par}\addvspace{3pt}\noindent\ignorespaces}
% 説明の段の見出し（丸数字＋太字の1行）  \bkivestep{番号}{見出し}
\newcommand\bkivestep[2]{\par\addvspace{6pt}\Needspace{5\baselineskip}\noindent
  {\bkcircled[bkOrange]{#1}\enspace\sffamily\bfseries\color{bkOrange!75!black}#2}\par\nobreak\vspace{2pt}\noindent\ignorespaces}
% 枠でかこんだ大事な式  \bkivewaku{…}
\newcommand\bkivewaku[1]{\par\addvspace{4pt}{\centering\setlength\fboxsep{5pt}\fcolorbox{bkOrange}{white}{#1}\par}\addvspace{4pt}\noindent\ignorespaces}
% 直角のしるし  \bkiveright{(頂点)}{(一方の点)}{(もう一方の点)}
\newcommand\bkiveright[3]{%
  \coordinate (bkivep) at ($#1!2.2mm!#2$);\coordinate (bkiveq) at ($#1!2.2mm!#3$);
  \draw[bkInk, line width=0.5pt] (bkivep) -- ($(bkivep)+(bkiveq)-#1$) -- (bkiveq);}
% 方眼の上のかたむいた正方形  \bkivetilt{(左下)}{a}{b}{1マスの mm}
%   外の正方形（1辺 a+b）、すみの直角三角形4つ（灰色）、かたむいた正方形（青）
\newcommand\bkivetilt[4]{\begin{scope}[shift={#1}, x=#4mm, y=#4mm]
  \pgfmathsetmacro\bkiveS{#2+#3}%
  \fill[bkBlue!18] (#3,0) -- (\bkiveS,#3) -- (#2,\bkiveS) -- (0,#2) -- cycle;
  \fill[bkGray!22] (0,0) -- (#3,0) -- (0,#2) -- cycle;
  \fill[bkGray!22] (#3,0) -- (\bkiveS,0) -- (\bkiveS,#3) -- cycle;
  \fill[bkGray!22] (\bkiveS,#3) -- (\bkiveS,\bkiveS) -- (#2,\bkiveS) -- cycle;
  \fill[bkGray!22] (#2,\bkiveS) -- (0,\bkiveS) -- (0,#2) -- cycle;
  \draw[bkLine, line width=0.3pt, step=1] (0,0) grid (\bkiveS,\bkiveS);
  \draw[bkInk!70, line width=0.6pt] (0,0) rectangle (\bkiveS,\bkiveS);
  \draw[bkBlue!80!black, line width=1.1pt] (#3,0) -- (\bkiveS,#3) -- (#2,\bkiveS) -- (0,#2) -- cycle;
\end{scope}}
\makeatother

\bkchapter{4}{4}{三平方の定理}{\bklevel{3}}{三平方の定理}
\bkmichishirube{2-8,3-2,3-8,3-10,4-1,4-3}{4-6,4-8,4-10,4-13,5-16}

\begin{goal}
\begin{itemize}[label={\color{\bkTT}●}, leftmargin=1.6\zw, itemsep=0pt, topsep=1pt]
\item \bklevel{3}\enspace 方眼の正方形の面積から、直角三角形の3つの辺の関係 $a^2+b^2=c^2$ を見つけ、2通りの方法でわけを説明できる。
\item \bklevel{3}\enspace 三平方の定理を使って、辺の長さ・2点の距離・直方体の対角線を求められる。
\item \bklevel{3}\enspace 3辺の長さから直角三角形かどうかを見分けられ、円の式 $x^2+y^2=r^2$ のわけがわかる。
\end{itemize}
\end{goal}

\begin{hatena}
\noindent\begin{minipage}[c]{0.57\linewidth}
3-2では、方眼の上の3つの点 A$(-2,\ 1)$、B$(3,\ 1)$、C$(3,\ -2)$ で、横に並んだAとBの距離が5、たてに並んだBとCの距離が3だと求めました。

では、ななめに並んだAとCの距離はいくつでしょう。Bを通って行くと $5+3=8$ ですが、まっすぐ行くほうが近道なので、8より短いはずです。ものさしではかると、5.8くらい。

たてと横の長さだけから、ななめの長さを計算で出す方法はあるでしょうか。
\end{minipage}\hfill
\begin{minipage}[c]{0.4\linewidth}\centering
\begin{tikzpicture}[x=5.4mm, y=5.4mm]
  \draw[bkLine, line width=0.3pt, step=1] (-3,-3) grid (4,2);
  \draw[bkGray, line width=0.6pt, -{Stealth[length=4pt]}] (-3.3,0) -- (4.5,0) node[anchor=west, font=\sffamily\footnotesize] {$x$};
  \draw[bkGray, line width=0.6pt, -{Stealth[length=4pt]}] (0,-3.3) -- (0,2.6) node[anchor=south, font=\sffamily\footnotesize] {$y$};
  \draw[bkBlue, line width=1.3pt] (-2,1) -- (3,1);
  \draw[bkOrange, line width=1.3pt] (3,1) -- (3,-2);
  \draw[bkRed, line width=1.3pt, dashed] (-2,1) -- (3,-2);
  \fill[bkInk] (-2,1) circle (0.13) (3,1) circle (0.13) (3,-2) circle (0.13);
  \node[font=\sffamily\footnotesize, anchor=south east, fill=white, inner sep=1pt] at (-2.1,1.15) {A};
  \node[font=\sffamily\footnotesize, anchor=south west, fill=white, inner sep=1pt] at (3.1,1.15) {B};
  \node[font=\sffamily\footnotesize, anchor=west, fill=white, inner sep=1pt] at (3.2,-2) {C};
  \node[font=\sffamily\footnotesize, text=bkBlue!80!black, anchor=south, fill=white, inner sep=1pt] at (1.5,1.25) {5};
  \node[font=\sffamily\footnotesize, text=bkOrange!80!black, anchor=west, fill=white, inner sep=1pt] at (3.25,-0.5) {3};
  \node[font=\sffamily\bfseries\small, text=bkRed!80!black, fill=white, inner sep=1pt] at (0.2,-1.3) {？};
\end{tikzpicture}
\end{minipage}
\end{hatena}

\noindent ななめの距離は、いろいろな所で必要になります。地図の上の2つの場所の近さ、テレビの画面の大きさ、ロボットの手が届くかどうか、地震の震源までの距離などです。この章では、直角三角形の3つの辺のあいだに成り立つ、とてもすっきりしたきまりを、方眼の上の正方形の面積から見つけます。そして、そのきまりがどんな直角三角形でも成り立つわけを、2通りの方法で確かめます。

%% ---------------- やってみよう・思い出そう ----------------
\begin{omoidasou}
\bkrecall{2-8}{2乗すると□になる正の数を $\sqrt{\text{□}}$ と書く。1辺1の正方形の対角線は $\sqrt2$。}
\bkrecall{3-2}{横に並んだ2点の距離は $|x_2-x_1|$、たてに並んだ2点の距離は $|y_2-y_1|$。}
\bkrecall{3-8}{$(a+b)^2=a^2+2ab+b^2$（1辺 $a+b$ の正方形を4つに分けた面積の図）。}
\bkrecall{3-10}{$\sqrt{12}=\sqrt{4\times3}=2\sqrt3$（$\sqrt{\phantom{2}}$ の中はできるだけ小さい整数に）。}
\bkrecall{4-1}{定理は、証明されたきまり。「AならばB」の逆は「BならばA」。3組の辺が等しい三角形は合同。三角形の内角の和は180°。}
\bkrecall{4-3}{2組の角がそれぞれ等しい三角形は相似で、対応する辺の比が等しい。}
\end{omoidasou}

\begin{yattemiyou}[方眼に、かたむいた正方形をかこう]
\noindent\begin{minipage}[c]{0.62\linewidth}
\begin{enumerate}[label={\bkcircled[bkGreen]{\arabic*}}, leftmargin=2.4\zw, itemsep=1pt, topsep=0pt]
\item 方眼紙のマスの角に、点Aをうつ。
\item Aから「右へ2、上へ1」の所に点B、Bから「左へ1、上へ2」の所に点C、Cから「左へ2、下へ1」の所に点Dをうつ。
\item A、B、C、Dを順に結ぶと、かたむいた正方形ができる（Dから「右へ1、下へ2」でAにもどる）。
\item この正方形の面積は何マス分かな？\quad 数えにくいときは、外側の大きな正方形から、すみの4つの直角三角形（灰色）をひこう。
\end{enumerate}
\end{minipage}\hfill
\begin{minipage}[c]{0.34\linewidth}\centering
\begin{tikzpicture}
  \bkivetilt{(0,0)}{2}{1}{9}
  \begin{scope}[x=9mm, y=9mm]
  \node[font=\sffamily\footnotesize, anchor=north, fill=bkGreenBg, inner sep=1pt] at (1,-0.15) {A};
  \node[font=\sffamily\footnotesize, anchor=west, fill=bkGreenBg, inner sep=1pt] at (3.15,1) {B};
  \node[font=\sffamily\footnotesize, anchor=south, fill=bkGreenBg, inner sep=1pt] at (2,3.15) {C};
  \node[font=\sffamily\footnotesize, anchor=east, fill=bkGreenBg, inner sep=1pt] at (-0.15,2) {D};
  \end{scope}
\end{tikzpicture}
\end{minipage}
\par\vspace{4pt}
\begin{enumerate}[label={\bkcircled[bkGreen]{\arabic*}}, leftmargin=2.4\zw, itemsep=1pt, topsep=0pt, start=5]
\item 「右へ3、上へ1」「右へ3、上へ2」でも、同じように正方形をかいて、面積を求めよう。
\end{enumerate}
\begin{bkcenter}
\begin{bktable}{colspec={Q[c,m,wd=30mm] Q[c,m,wd=22mm] Q[c,m,wd=28mm] Q[c,m,wd=28mm]}, rowsep=1pt, row{1}={font=\sffamily\footnotesize}}
はじめの1歩 & 外側の正方形 & すみの三角形4つ & かたむいた正方形 \\
右へ2、上へ1 & 9マス & & \\
右へ3、上へ1 & & & \\
右へ3、上へ2 & & & \\
\end{bktable}
\end{bkcenter}
\end{yattemiyou}

\noindent やってみようの答え：「右へ2、上へ1」では、外側の正方形は1辺が $2+1=3$ マスなので、$3\times3=9$ マス。すみの直角三角形は、たて1・横2の長方形の半分なので $2\times1\div2=1$ マス。それが4つで4マス。だから、かたむいた正方形は $9-4=5$ マスです。同じように計算すると、「右へ3、上へ1」は $16-4\times1.5=10$ マス、「右へ3、上へ2」は $25-4\times3=13$ マスになります。


%% ---------------- 数学の道具 ----------------
\bksection{doguu}{数学の道具}

\bksubsection{1}{方眼の正方形で、ななめの長さを調べる}
\label{bk:4-4:mitsukeru}やってみようで作った正方形を、もう一度ながめます。かたむいた正方形の1辺は、方眼の上を「右へ $a$、上へ $b$」進んだ線分です。この線分は、たて $b$・横 $a$ の直角三角形の、ななめの辺になっています。

\begin{bkfigure}
\begin{tikzpicture}[x=1mm, y=1mm]
  \bkivetilt{(0,0)}{3}{1}{6}
  \bkivetilt{(36,0)}{3}{2}{6}
  \bkivetilt{(78,0)}{4}{3}{6}
  \node[font=\sffamily\footnotesize, anchor=north, align=center] at (12,-2) {右へ3、上へ1\\$16-4\times1.5=10$};
  \node[font=\sffamily\footnotesize, anchor=north, align=center] at (51,-2) {右へ3、上へ2\\$25-4\times3=13$};
  \node[font=\sffamily\footnotesize, anchor=north, align=center] at (99,-2) {右へ4、上へ3\\$49-4\times6=25$};
\end{tikzpicture}
\bkfigcaption{外側の正方形から、すみの直角三角形（灰色）4つをひくと、かたむいた正方形（青）の面積になる。}
\end{bkfigure}

いろいろな「はじめの1歩」で正方形をかき、面積を表にまとめます。$a^2$ と $b^2$ の列も、となりにならべてみましょう。

\begin{bkcenter}
\begin{bktable}{colspec={Q[c,m,wd=22mm] Q[c,m,wd=22mm] Q[c,m,wd=16mm] Q[c,m,wd=16mm] Q[c,m,wd=34mm]}, rowsep=1.5pt, row{1}={font=\sffamily\small}}
右へ $a$ & 上へ $b$ & $a^2$ & $b^2$ & かたむいた正方形の面積 \\
1 & 1 & 1 & 1 & 2 \\
2 & 1 & 4 & 1 & 5 \\
3 & 1 & 9 & 1 & 10 \\
3 & 2 & 9 & 4 & 13 \\
4 & 3 & 16 & 9 & 25 \\
\end{bktable}
\bktabcaption{どの行も、$a^2$ と $b^2$ をたすと、かたむいた正方形の面積になっている。}
\end{bkcenter}

どの行でも、$a^2+b^2$ が、かたむいた正方形の面積と同じになっています。$1+1=2$、$4+1=5$、$9+1=10$、$9+4=13$、$16+9=25$ です。

\bkivehead{面積がわかれば、1辺がわかる}
正方形の面積は「1辺×1辺」なので、1辺は「2乗すると面積になる正の数」、つまり $\sqrt{\text{面積}}$ です（2-8）。
\begin{itemize}[leftmargin=1.2\zw, itemsep=1pt, topsep=2pt, label={\color{\bkTT}●}]
\item 右へ4、上へ3：面積25なので、1辺は $\sqrt{25}=5$。ものさしではかっても5マス分です。
\item 右へ2、上へ1：面積5なので、1辺は $\sqrt5=2.236\cdots$。ものさしではかると2.2マス分くらいで、合っています。
\item 右へ1、上へ1：面積2なので、1辺は $\sqrt2$。2-8で見た「正方形の対角線が $\sqrt2$」と同じです。
\end{itemize}
こうして、ななめの線分の長さが、たてと横の長さから計算できました。

\bkivehead{直角三角形の3つの辺に、正方形をのせる}
方眼の上の直角三角形（直角をはさむ辺が3と4）の、3つの辺のそれぞれに、正方形をのせてみます。

\noindent\begin{minipage}[c]{0.5\linewidth}\centering
\begin{tikzpicture}[x=4.4mm, y=4.4mm]
  % 直角三角形：C(0,0)、B(4,0)、A(0,3)
  \fill[bkOrange!22] (0,0) rectangle (4,-4);
  \draw[bkLine, line width=0.3pt, step=1] (0,-4) grid (4,0);
  \draw[bkOrange!80!black, line width=0.8pt] (0,0) rectangle (4,-4);
  \fill[bkBlue!18] (0,0) rectangle (-3,3);
  \draw[bkLine, line width=0.3pt, step=1] (-3,0) grid (0,3);
  \draw[bkBlue!80!black, line width=0.8pt] (0,0) rectangle (-3,3);
  \fill[bkGreen!20] (0,3) -- (4,0) -- (7,4) -- (3,7) -- cycle;
  \foreach \k in {1,2,3,4} {
    \draw[bkLine, line width=0.3pt] ($(0,3)+\k*(0.8,-0.6)$) -- ++(3,4);
    \draw[bkLine, line width=0.3pt] ($(0,3)+\k*(0.6,0.8)$) -- ++(4,-3);}
  \draw[bkGreen!60!black, line width=0.8pt] (0,3) -- (4,0) -- (7,4) -- (3,7) -- cycle;
  \fill[bkYellow!35] (0,0) -- (4,0) -- (0,3) -- cycle;
  \draw[bkInk, line width=1.1pt] (0,0) -- (4,0) -- (0,3) -- cycle;
  \bkiveright{(0,0)}{(4,0)}{(0,3)}
  \node[font=\sffamily\small, fill=white, inner sep=1.5pt] at (2,-2) {16};
  \node[font=\sffamily\small, fill=white, inner sep=1.5pt] at (-1.5,1.5) {9};
  \node[font=\sffamily\small, fill=white, inner sep=1.5pt] at (3.5,3.5) {25};
\end{tikzpicture}
\end{minipage}\hfill
\begin{minipage}[c]{0.47\linewidth}
直角をはさむ辺にのせた正方形は、9マスと16マスです。いちばん長い辺にのせた、かたむいた正方形は、表の「右へ4、上へ3」と同じで25マスです。
\bkiveline{$9+16=25$}
小さい2つの正方形の面積をたすと、大きい正方形の面積とぴったり同じになります。
\end{minipage}

\bkivehead{直角でないと、どうなる？}
2つの辺を \qty{3}{cm} と \qty{4}{cm} にしたまま、そのあいだの角を変えてみます。3つめの辺の長さをはかり、2乗してみましょう。

\begin{bkfigure}
\begin{tikzpicture}[x=1mm, y=1mm, lb/.style={font=\sffamily\footnotesize, anchor=north, align=center}]
  % 1 cm を 6 mm でかく
  \begin{scope}[shift={(4,0)}]
    \fill[bkBlue!10] (0,0) -- (24,0) -- (9,15.59) -- cycle;
    \draw[bkInk, line width=0.9pt] (0,0) -- (24,0) (0,0) -- (9,15.59);
    \draw[bkRed, line width=1.3pt] (24,0) -- (9,15.59);
    \draw[bkInk, line width=0.5pt] (5,0) arc (0:60:5);
    \node[font=\sffamily\footnotesize] at (30:8) {60°};
    \node[font=\sffamily\footnotesize, anchor=east] at (2.8,8.5) {\qty{3}{cm}};
    \node[lb] at (12,-7) {約\qty{3.61}{cm}\\2乗すると13};
  \end{scope}
  \begin{scope}[shift={(48,0)}]
    \fill[bkGreen!12] (0,0) -- (24,0) -- (0,18) -- cycle;
    \draw[bkInk, line width=0.9pt] (0,0) -- (24,0) (0,0) -- (0,18);
    \draw[bkRed, line width=1.3pt] (24,0) -- (0,18);
    \bkiveright{(0,0)}{(24,0)}{(0,18)}
    \node[font=\sffamily\footnotesize, anchor=east] at (-1,9) {\qty{3}{cm}};
    \node[lb] at (12,-7) {\qty{5}{cm}\\2乗すると25};
  \end{scope}
  \begin{scope}[shift={(96,0)}]
    \fill[bkOrange!12] (0,0) -- (24,0) -- (-9,15.59) -- cycle;
    \draw[bkInk, line width=0.9pt] (0,0) -- (24,0) (0,0) -- (-9,15.59);
    \draw[bkRed, line width=1.3pt] (24,0) -- (-9,15.59);
    \draw[bkInk, line width=0.5pt] (4,0) arc (0:120:4);
    \node[font=\sffamily\footnotesize] at (60:7.5) {120°};
    \node[font=\sffamily\footnotesize, anchor=east] at (-6.2,7.5) {\qty{3}{cm}};
    \node[lb] at (8,-7) {約\qty{6.08}{cm}\\2乗すると37};
  \end{scope}
  \foreach \s in {4,48,96} {\node[font=\sffamily\footnotesize, anchor=north] at (\s+12,-1) {\qty{4}{cm}};}
\end{tikzpicture}
\bkfigcaption{2辺は \qty{3}{cm} と \qty{4}{cm}。間の角が60°・90°・120°のときの、3つめの辺（赤）。}
\end{bkfigure}

$3^2+4^2=25$ になるのは、まん中の直角のときだけです。角が直角より小さいと、3つめの辺は短くなり、2乗は25より小さくなります（13）。角が直角より大きいと、長くなり、2乗は25より大きくなります（37）。\textbf{$a^2+b^2$ が3つめの辺の2乗と同じになるのは、直角三角形のときのきまり}です。

\bkivehead{このきまりに、名前をつける}
直角三角形で、直角の向かいにある辺を\textbf{斜辺（しゃへん）}といいます。直角をはさむ2辺の長さを $a$、$b$、斜辺の長さを $c$ とすると、表や図で見たことは、次のように書けます。\bkidx{しゃへん}{斜辺}
\bkivewaku{$a^2+b^2=c^2$}
これを\textbf{三平方の定理（さんへいほうのていり）}といいます。3つの辺の平方（2乗）のあいだのきまりなので「三平方」です。2500年ほど前のギリシャのピタゴラスの名前をとって、\textbf{ピタゴラスの定理}ともいいます。\bkidx{さんへいほうのていり}{三平方の定理}

$c^2$ は $a^2$ に $b^2$ をたした数なので、$a^2$ よりも $b^2$ よりも大きくなります。だから、\textbf{斜辺は、3つの辺の中でいちばん長い辺}です。

ただし、ここまでは、いくつかの直角三角形で確かめただけです。4-1で学んだように、何回確かめても、それだけでは「どんな直角三角形でも必ず」とは言えません。「定理」とよぶためには、証明が必要です。

\bksubsection{2}{どんな直角三角形でも成り立つわけ：2通りの証明}
\label{bk:4-4:shoumei}ここでは、直角三角形の辺の長さを、数ではなく文字 $a$、$b$、$c$ のままにして考えます。文字のまま説明できれば、どんな大きさ・どんな形の直角三角形にもあてはまります。

\bkivestep{1}{証明1：同じ板に、同じ三角形を4まいずつ置く}
1辺が $a+b$ の正方形の板を2まい用意します。そして、直角をはさむ辺が $a$ と $b$、斜辺が $c$ の直角三角形（灰色）を、それぞれの板に4まいずつ置きます。置き方だけを変えます。

\begin{bkfigure}
\begin{tikzpicture}[x=1mm, y=1mm]
  % 左の板：すみに三角形4まい、まん中に c の正方形（a=24、b=12）
  \fill[bkGreen!22] (24,0) -- (36,24) -- (12,36) -- (0,12) -- cycle;
  \fill[bkGray!28] (0,0) -- (24,0) -- (0,12) -- cycle;
  \fill[bkGray!28] (24,0) -- (36,0) -- (36,24) -- cycle;
  \fill[bkGray!28] (36,24) -- (36,36) -- (12,36) -- cycle;
  \fill[bkGray!28] (12,36) -- (0,36) -- (0,12) -- cycle;
  \draw[bkInk, line width=0.6pt] (24,0) -- (36,24) -- (12,36) -- (0,12) -- cycle;
  \draw[bkInk, line width=0.9pt] (0,0) rectangle (36,36);
  \node[font=\small] at (18,18) {$c^2$};
  \node[font=\small] at (13.6,9.2) {$c$};
  \bkbrace{(0,-1.5)}{(24,-1.5)}{$a$}{-1}\bkbrace{(24,-1.5)}{(36,-1.5)}{$b$}{-1}
  \node[font=\small, anchor=east] at (-1.5,6) {$b$};
  \node[font=\small, anchor=east] at (-1.5,24) {$a$};
  \draw[bkInk!60, line width=0.4pt] (-1,12) -- (0,12);
  % 右の板：三角形2まいずつで長方形2つ。のこりは a² と b²
  \begin{scope}[shift={(66,0)}]
  \fill[bkBlue!20] (0,12) rectangle (24,36);
  \fill[bkOrange!25] (24,0) rectangle (36,12);
  \fill[bkGray!28] (0,0) rectangle (24,12);
  \fill[bkGray!28] (24,12) rectangle (36,36);
  \draw[bkInk, line width=0.6pt] (0,12) -- (24,0) (24,36) -- (36,12);
  \draw[bkInk, line width=0.6pt] (0,12) -- (24,12) -- (24,0) (24,12) -- (36,12) (24,12) -- (24,36);
  \draw[bkInk, line width=0.9pt] (0,0) rectangle (36,36);
  \node[font=\small] at (12,24) {$a^2$};
  \node[font=\small] at (30,6) {$b^2$};
  \bkbrace{(0,-1.5)}{(24,-1.5)}{$a$}{-1}\bkbrace{(24,-1.5)}{(36,-1.5)}{$b$}{-1}
  \node[font=\small, anchor=west] at (37.5,6) {$b$};
  \node[font=\small, anchor=west] at (37.5,24) {$a$};
  \draw[bkInk!60, line width=0.4pt] (36,12) -- (37,12);
  \end{scope}
  \node[font=\sffamily\footnotesize, anchor=north] at (18,-9) {左の置き方};
  \node[font=\sffamily\footnotesize, anchor=north] at (84,-9) {右の置き方};
\end{tikzpicture}
\bkfigcaption{同じ大きさの板に、同じ直角三角形（灰色）を4まいずつ置く。三角形にかくれていない部分は、左は $c^2$、右は $a^2$ と $b^2$。}\label{bk:4-4:fig-ita}
\end{bkfigure}

\begin{enumerate}[label={\bkcircled[bkOrange]{\arabic*}}, leftmargin=2.4\zw, itemsep=2pt, topsep=2pt]
\item \textbf{左の置き方}：板の4つのすみに三角形を置くと、まん中に、4つの辺がどれも $c$ の四角形がのこります。この四角形の角は、どれも直角です。なぜなら、板の辺の上の1点に集まる角は「三角形の直角でない2つの角」と「四角形の角」で、合わせて180°。三角形の直角でない2つの角の和は $180^\circ-90^\circ=90^\circ$（内角の和、4-1）なので、四角形の角は $180^\circ-90^\circ=90^\circ$ です。だから、のこった部分は1辺 $c$ の正方形で、面積は $c^2$ です。
\item \textbf{右の置き方}：三角形を2まいずつ組にして、2辺が $a$ と $b$ の長方形を2つ作ります。すると、のこった部分は、1辺 $a$ の正方形（面積 $a^2$）と、1辺 $b$ の正方形（面積 $b^2$）です。
\item 板も、置いた三角形4まいも、左右で同じです。だから、三角形にかくれていない部分の面積も同じです。
\end{enumerate}
\bkiveline{$a^2+b^2=c^2$}

\bkivehead{式でも確かめる（3-8の展開）}
左の板の面積を、2通りに表します。板は1辺 $a+b$ の正方形なので、面積は $(a+b)^2$。いっぽう、まん中の正方形 $c^2$ と、三角形4まい（1まい $\dfrac{ab}{2}$）を合わせた面積でもあります。
\begin{align*}
(a+b)^2&=c^2+4\times\frac{ab}{2}&&\text{（板の面積を2通りに表した）}\\
a^2+2ab+b^2&=c^2+2ab&&\text{（左は展開、3-8。右は $4\times\dfrac{ab}{2}=2ab$）}\\
a^2+b^2&=c^2&&\text{（両辺から $2ab$ をひいた）}
\end{align*}
図で考えても、式で考えても、同じ結果になりました。

\bkfukidashi{mina}{証明を1つ読みきったね。ここまで来たら、もう1つ、ちがう見方の証明も見てみよう。}

\bkivestep{2}{証明2：相似な三角形を見つける（4-3）}
直角三角形ABC（∠C＝90°）で、直角の頂点Cから、斜辺ABに垂直な線をひき、ABと交わる点をHとします。すると、大きな直角三角形の中に、小さな直角三角形が2つ（青とオレンジ）できます。

\noindent\begin{minipage}[c]{0.46\linewidth}\centering
\begin{tikzpicture}[x=1mm, y=1mm]
  \coordinate (A) at (0,0); \coordinate (B) at (50,0); \coordinate (C) at (18,24); \coordinate (H) at (18,0);
  \fill[bkBlue!16] (A) -- (H) -- (C) -- cycle;
  \fill[bkOrange!20] (H) -- (B) -- (C) -- cycle;
  \draw[bkInk, line width=0.6pt] (C) -- (H);
  \draw[bkInk, line width=1.0pt] (A) -- (B) -- (C) -- cycle;
  \bkiveright{(C)}{(A)}{(B)}
  \bkiveright{(H)}{(B)}{(C)}
  \draw[bkBlue!80!black, line width=0.9pt] (6,0) arc (0:53.13:6);
  \draw[bkOrange!80!black, line width=0.9pt] (44,0) arc (180:143.13:6);
  \draw[bkOrange!80!black, line width=0.9pt] ($(C)+(-126.87:8)$) arc (-126.87:-90:8);
  \draw[bkBlue!80!black, line width=0.9pt] ($(C)+(-90:6.5)$) arc (-90:-36.87:6.5);
  \node[font=\sffamily\small, anchor=east] at (-1,0) {A};
  \node[font=\sffamily\small, anchor=west] at (51,0) {B};
  \node[font=\sffamily\small, anchor=south] at (18,25.2) {C};
  \node[font=\sffamily\small, anchor=north] at (18,-1) {H};
  \node[font=\small, anchor=south east] at (8.4,12.4) {$b$};
  \node[font=\small, anchor=south west] at (34.6,12.6) {$a$};
  \bkbrace{(0,-7)}{(50,-7)}{$c$}{-1}
\end{tikzpicture}
\bkfigcaption{同じ色の角は、同じ大きさ。}\label{bk:4-4:fig-souji}
\end{minipage}\hfill
\begin{minipage}[c]{0.51\linewidth}
大きな三角形ACBと、青い三角形AHCをくらべます。
\begin{itemize}[leftmargin=1.2\zw, itemsep=1pt, topsep=2pt, label={\color{\bkTT}●}]
\item ∠Aは、どちらにも共通（青い角）。
\item ∠ACB＝∠AHC＝90°。
\end{itemize}
2組の角がそれぞれ等しいので、2つの三角形は相似です（4-3）。のこりの角（オレンジ）も等しくなります。
\end{minipage}

\vspace{3pt}
相似な三角形では、対応する辺の比が等しいので、「斜辺：Aをはさむ短いほうの辺」の比をくらべます。大きな三角形では AB：AC、青い三角形では AC：AH です。
\begin{align*}
\frac{\mathrm{AB}}{\mathrm{AC}}&=\frac{\mathrm{AC}}{\mathrm{AH}}&&\text{（対応する辺の比が等しい）}\\
\mathrm{AB}\times\mathrm{AH}&=\mathrm{AC}\times\mathrm{AC}&&\text{（両辺に $\mathrm{AC}\times\mathrm{AH}$ をかけた、2-14）}\\
c\times\mathrm{AH}&=b^2
\end{align*}
オレンジの三角形BHCと大きな三角形BCAでも、∠Bが共通で直角があるので相似です。同じように考えると、$c\times\mathrm{BH}=a^2$ です。2つの式をたします。
\begin{align*}
a^2+b^2&=c\times\mathrm{BH}+c\times\mathrm{AH}\\
&=c\times(\mathrm{BH}+\mathrm{AH})&&\text{（$c$ でくくった、3-8）}\\
&=c\times c=c^2&&\text{（BHとAHを合わせると、斜辺ABの長さ $c$）}
\end{align*}

\bkivehead{数で確かめる}
3辺が \qty{3}{cm}・\qty{4}{cm}・\qty{5}{cm} の直角三角形（$b=3$、$a=4$、$c=5$）で確かめます。$c\times\mathrm{AH}=b^2$ から $\mathrm{AH}=9\div5=\qty{1.8}{cm}$、$c\times\mathrm{BH}=a^2$ から $\mathrm{BH}=16\div5=\qty{3.2}{cm}$ です。たすと $1.8+3.2=5$ で、ちゃんと斜辺の長さになります。

\bkivehead{証明ができたので「定理」}
2つの証明は、どちらも文字 $a$、$b$、$c$ のままで考えたので、どんな直角三角形にもあてはまります。何回はかっても言えなかった「必ず」が、これで言えるようになりました。証明されたきまりなので、「三平方の\textbf{定理}」とよびます（定義・定理・法則のちがいは4-1）。

\Needspace{8\baselineskip}
\begin{seikaku}
三平方の定理は、平らな面の上にかいた直角三角形のきまりです。地球の表面のように曲がった面の上に、とても大きな三角形をかくと、成り立ちません。たとえば、北極と、赤道の上で経度が90°ちがう2つの点を結ぶ三角形は、3つの角がすべて90°で、3つの辺がすべて約 \qty{10000}{km} です。これでは $a^2+b^2=c^2$ になりません。内角の和が180°になる（4-1）のも、平らな面のきまりです。ふつうの大きさの図形や機械を考えるときは、平らな面のきまりで十分です。
\end{seikaku}

\bksubsection{3}{2つの辺から、のこりの辺を求める}
\label{bk:4-4:motomeru}直角三角形の2つの辺の長さがわかれば、三平方の定理で、のこりの辺の長さを計算できます。わかっている辺によって、2つの場合があります。

\bkivehead{はじめに、斜辺はどれかを確かめる}
斜辺は「直角の向かいにある辺」で、いちばん長い辺です。$a^2+b^2=c^2$ の $c$ は、いつも斜辺です。どの辺が斜辺かをまちがえると、答えがまったくちがってしまいます。

\noindent\begin{minipage}[c]{0.49\linewidth}
\bkivehead{(ア) 直角をはさむ2辺がわかる}
$a=6$、$b=8$（cm）のとき、
\begin{align*}
c^2&=6^2+8^2=36+64=100\\
c&=\sqrt{100}=10
\end{align*}
斜辺は \qty{10}{cm}。$c$ は、2乗すると100になる正の数です（2-8）。
\end{minipage}\hfill
\begin{minipage}[c]{0.49\linewidth}
\bkivehead{(イ) 斜辺と、もう1つの辺がわかる}
$c=10$、$b=6$（cm）のとき、$a^2+b^2=c^2$ を $a^2$ について解いて（2-14）、
\begin{align*}
a^2&=c^2-b^2=100-36=64\\
a&=\sqrt{64}=8
\end{align*}
のこりの辺は \qty{8}{cm}。
\end{minipage}

\vspace{4pt}
2乗して100になる数には、$-10$ もあります（$(-10)^2=100$、2-8）。でも、長さはマイナスにならないので、正のほうの10だけが答えです。

\bkivehead{答えが整数にならないとき}
直角をはさむ辺が \qty{1}{cm} と \qty{2}{cm} のときは、
\bkiveline{$c^2=1^2+2^2=5$\qquad $c=\sqrt5=2.236\cdots$}
なので、約 \qty{2.24}{cm} です（$\sqrt5$ の値は電卓で求める）。$\sqrt{\phantom{2}}$ のまま答えるときは、中をできるだけ小さい整数にします（3-10）。2辺が2と4なら、
\bkiveline{$c=\sqrt{2^2+4^2}=\sqrt{20}=\sqrt{4\times5}=2\sqrt5$}

\bkivehead{3辺がどれも整数になる直角三角形}
3辺がどれも整数になる組は、とくべつです。よく出てくるものを表にしておきます。
\begin{bkcenter}
\begin{bktable}{colspec={Q[c,m,wd=26mm] Q[c,m,wd=26mm] Q[c,m,wd=20mm] Q[c,m,wd=34mm]}, rowsep=1.5pt, row{1}={font=\sffamily\small}}
\SetCell[c=2]{c} 直角をはさむ辺 & & 斜辺 & 確かめ \\
3 & 4 & 5 & $9+16=25$ \\
5 & 12 & 13 & $25+144=169$ \\
8 & 15 & 17 & $64+225=289$ \\
7 & 24 & 25 & $49+576=625$ \\
6 & 8 & 10 & $36+64=100$ \\
\end{bktable}
\bktabcaption{6・8・10は、3・4・5を2倍にした相似な三角形（4-3）。}
\end{bkcenter}

\bksubsection{4}{三角定規の直角三角形：$1:1:\sqrt2$ と $1:2:\sqrt3$}
\label{bk:4-4:jougi}三角定規は、2まいで1組です。1まいは直角二等辺三角形（角が45°・45°・90°）、もう1まいは角が30°・60°・90°の直角三角形です。この2つの辺の比を、三平方の定理で求めます。

\begin{bkfigure}
\begin{tikzpicture}[x=1mm, y=1mm]
  % 左：1辺1の正方形と対角線（1を28mm）
  \fill[bkBlue!16] (0,0) -- (28,0) -- (28,28) -- cycle;
  \draw[bkInk!60, line width=0.6pt, dashed] (0,0) -- (0,28) -- (28,28);
  \draw[bkInk, line width=1.0pt] (0,0) -- (28,0) -- (28,28) -- cycle;
  \bkiveright{(28,0)}{(0,0)}{(28,28)}
  \draw[bkInk, line width=0.5pt] (6,0) arc (0:45:6);
  \node[font=\sffamily\footnotesize] at (22.5:9.6) {45°};
  \draw[bkInk, line width=0.5pt] ($(28,28)+(-90:6)$) arc (-90:-135:6);
  \node[font=\sffamily\footnotesize] at ($(28,28)+(-112.5:9.6)$) {45°};
  \node[font=\small, anchor=north] at (14,-1.5) {1};
  \node[font=\small, anchor=west] at (29.5,14) {1};
  \node[font=\small, anchor=south east] at (13,15) {$\sqrt2$};
  \node[font=\sffamily\footnotesize, anchor=north, align=center] at (14,-8) {直角二等辺三角形\\$1:1:\sqrt2$};
  % 右：1辺2の正三角形を半分に（1を20mm）
  \begin{scope}[shift={(66,0)}]
  \fill[bkOrange!20] (20,0) -- (40,0) -- (20,34.64) -- cycle;
  \draw[bkInk!60, line width=0.6pt, dashed] (0,0) -- (20,0) (0,0) -- (20,34.64);
  \draw[bkInk, line width=1.0pt] (20,0) -- (40,0) -- (20,34.64) -- cycle;
  \bkiveright{(20,0)}{(40,0)}{(20,34.64)}
  \draw[bkInk, line width=0.5pt] ($(40,0)+(180:6)$) arc (180:120:6);
  \node[font=\sffamily\footnotesize] at ($(40,0)+(150:9.4)$) {60°};
  \draw[bkInk, line width=0.5pt] ($(20,34.64)+(-90:7)$) arc (-90:-60:7);
  \node[font=\sffamily\footnotesize] at ($(20,34.64)+(-75:10.6)$) {30°};
  \node[font=\small, anchor=north] at (30,-1.5) {1};
  \node[font=\small, anchor=south west] at (30.5,17.5) {2};
  \node[font=\small, anchor=east] at (18.5,13) {$\sqrt3$};
  \node[font=\sffamily\footnotesize, anchor=north, align=center] at (20,-8) {正三角形の半分\\$1:2:\sqrt3$};
  \end{scope}
\end{tikzpicture}
\bkfigcaption{左：1辺1の正方形を対角線で半分に。右：1辺2の正三角形を、頂点からまっすぐ下ろした線で半分に。}\label{bk:4-4:fig-jougi}
\end{bkfigure}

\bkivehead{直角二等辺三角形：$1:1:\sqrt2$}
1辺1の正方形を、対角線で半分に切ります。正方形は対角線について線対称（4-1）なので、直角が半分ずつに分かれて45°になります。斜辺は
\bkiveline{$c^2=1^2+1^2=2$\qquad だから\qquad $c=\sqrt2$}
2-8では、面積2の正方形を作って「1辺1の正方形の対角線は $\sqrt2$」と知りました。三平方の定理でも、同じ答えになります。

\bkivehead{30°・60°・90°の直角三角形：$1:2:\sqrt3$}
1辺2の正三角形（3つの角はどれも60°）を、上の頂点からまっすぐ下ろした線で半分にします。底辺は半分の1、斜辺はもとの辺の2のままです。上の角は60°の半分の30°になります。高さを $h$ とすると、
\bkiveline{$h^2+1^2=2^2$\qquad $h^2=4-1=3$\qquad $h=\sqrt3$}
いちばん短い辺：斜辺：のこりの辺＝$1:2:\sqrt3$ です。\textbf{斜辺は2で、$\sqrt3$ ではありません}。$\sqrt3=1.732\cdots$ は2より小さいので、斜辺（いちばん長い辺）にはなりません。

\bkivehead{大きさが変わっても、比は同じ}
角が同じ直角三角形は、どれも相似です（4-3）。だから、大きさが変わっても、辺の比は $1:1:\sqrt2$、$1:2:\sqrt3$ のままです。たとえば、
\begin{itemize}[leftmargin=1.2\zw, itemsep=1pt, topsep=2pt, label={\color{\bkTT}●}]
\item 直角をはさむ辺が \qty{15}{cm} の45°の三角定規の斜辺は、$15\times\sqrt2=15\sqrt2$ で、約 \qty{21.2}{cm}。
\item 斜辺が \qty{20}{cm} の30°・60°の三角定規は、短い辺が $20\div2=\qty{10}{cm}$、のこりの辺が $10\sqrt3$ で約 \qty{17.3}{cm}。
\item 1辺 $a$ の正三角形の高さは、$1:2:\sqrt3$ の「2」が $a$ なので、$\dfrac{a}{2}\times\sqrt3=\dfrac{\sqrt3}{2}a$。
\end{itemize}
4-6では、この比が、坂の急さを表す sin・cos という数になります（これから学ぶしるしです。いまはわからなくてもだいじょうぶ）。

\bksubsection{5}{逆：3辺の長さから直角を作る}
\label{bk:4-4:gyaku}4-1で、「AならばB」に対して「BならばA」を\textbf{逆}ということを学びました。逆は、いつも正しいとはかぎりません。「正方形ならば長方形」は正しいけれど、逆の「長方形ならば正方形」は正しくありません（たて2・横3の長方形が反例）。

三平方の定理は「三角形が直角三角形（∠C＝90°）ならば、$a^2+b^2=c^2$」でした。その逆は、
\begin{taisetsu}
三角形の3辺の長さ $a$、$b$、$c$ が $a^2+b^2=c^2$ をみたすならば、\\
長さ $c$ の辺の向かいの角は直角である。
\end{taisetsu}
\noindent この逆は正しいことが、次のように説明できます。\bkidx{さんへいほうのていりのぎゃく}{三平方の定理の逆}
\begin{enumerate}[label={\bkcircled[bkOrange]{\arabic*}}, leftmargin=2.4\zw, itemsep=2pt, topsep=2pt]
\item $a^2+b^2=c^2$ をみたす三角形を、三角形Pとします。
\item 別に、直角をはさむ辺が $a$ と $b$ の直角三角形Qを作ります。Qの斜辺を $d$ とすると、三平方の定理から $d^2=a^2+b^2$ です。
\item ①と②から $d^2=c^2$。$d$ も $c$ も正の数なので、$d=c$ です。
\item PとQは、3組の辺がそれぞれ等しいので合同です（4-1）。合同な三角形の対応する角は等しいので、Pの角も直角です。
\end{enumerate}

\bkivehead{ひも1本で直角を作る}
ひもに、同じ間かくで12のしるしをつけて、輪にします。しるしの所を3つ選び、辺の長さが3と4と5の三角形になるように、ピンと張ります。すると、3の辺と4の辺のあいだに、直角ができます。$3^2+4^2=9+16=25=5^2$ だからです。

\noindent\begin{minipage}[c]{0.4\linewidth}\centering
\begin{tikzpicture}[x=7mm, y=7mm]
  \fill[bkYellow!25] (0,0) -- (4,0) -- (0,3) -- cycle;
  \draw[bkOrange!80!black, line width=1.4pt] (0,0) -- (4,0) -- (0,3) -- cycle;
  \bkiveright{(0,0)}{(4,0)}{(0,3)}
  \foreach \k in {0,1,2,3,4} {\fill[bkInk] (\k,0) circle (0.11);}
  \foreach \k in {1,2,3} {\fill[bkInk] (0,\k) circle (0.11);}
  \foreach \k in {1,2,3,4} {\fill[bkInk] ($(4,0)+\k*(-0.8,0.6)$) circle (0.11);}
  \node[font=\small, anchor=north] at (2,-0.3) {4};
  \node[font=\small, anchor=east] at (-0.3,1.5) {3};
  \node[font=\small, anchor=south west] at (2.1,1.6) {5};
  \node[font=\sffamily\footnotesize, anchor=north, align=center] at (2,-1.2) {しるしは全部で12こ};
\end{tikzpicture}
\end{minipage}\hfill
\begin{minipage}[c]{0.56\linewidth}
大昔のエジプトでも、こうして直角を作ったといわれます。今でも、建物の土台の線を引くときに、\qty{3}{m}・\qty{4}{m}・\qty{5}{m}（または \qty{0.9}{m}・\qty{1.2}{m}・\qty{1.5}{m}）で直角を確かめます。大工さんが使う大きなL字の道具「おおがね」も、3:4:5 で作られています。
\end{minipage}

\Needspace{14\baselineskip}
\begin{motto}[3]
\textbf{直角より小さいか、大きいか}\par
①で見たように、2辺が3と4でも、間の角が直角より小さいと3つめの辺の2乗は25より小さく（13）、直角より大きいと25より大きく（37）なりました。いちばん長い辺を $c$ として、$a^2+b^2$ と $c^2$ をくらべると、$c$ の向かいの角がわかります。
\begin{itemize}[leftmargin=1.2\zw, itemsep=0pt, topsep=2pt, label={\color{bkGray}●}]
\item $a^2+b^2=c^2$ なら、直角
\item $a^2+b^2>c^2$ なら、直角より小さい角（鋭角（えいかく））
\item $a^2+b^2<c^2$ なら、直角より大きい角（鈍角（どんかく））
\end{itemize}
どれだけ大きいか小さいかは、4-7の余弦定理で、角の大きさといっしょに計算できるようになります。
\end{motto}

\bksubsection{6}{座標の上の2点の距離}
\label{bk:4-4:kyori}3-2では、横に並んだ2点、たてに並んだ2点の距離を求めました。ななめに並んだ2点も、三平方の定理を使えば求められます。

\noindent\begin{minipage}[c]{0.5\linewidth}\centering
\begin{tikzpicture}[x=8mm, y=8mm]
  \draw[bkGray, line width=0.6pt, -{Stealth[length=4pt]}] (-0.4,0) -- (6.2,0) node[anchor=west, font=\sffamily\footnotesize] {$x$};
  \draw[bkGray, line width=0.6pt, -{Stealth[length=4pt]}] (0,-0.4) -- (0,5.2) node[anchor=south, font=\sffamily\footnotesize] {$y$};
  \node[font=\sffamily\footnotesize, anchor=north east] at (0,0) {O};
  \draw[bkGray, dashed, line width=0.4pt] (1,0) -- (1,1) (5,0) -- (5,1) (0,1) -- (1,1) (0,4) -- (5,4);
  \node[font=\footnotesize, anchor=north] at (1,-0.1) {$x_1$};
  \node[font=\footnotesize, anchor=north] at (5,-0.1) {$x_2$};
  \node[font=\footnotesize, anchor=east] at (-0.1,1) {$y_1$};
  \node[font=\footnotesize, anchor=east] at (-0.1,4) {$y_2$};
  \fill[bkYellow!25] (1,1) -- (5,1) -- (5,4) -- cycle;
  \draw[bkBlue, line width=1.3pt] (1,1) -- (5,1);
  \draw[bkOrange, line width=1.3pt] (5,1) -- (5,4);
  \draw[bkRed, line width=1.4pt] (1,1) -- (5,4);
  \bkiveright{(5,1)}{(1,1)}{(5,4)}
  \fill[bkInk] (1,1) circle (0.09) (5,4) circle (0.09) (5,1) circle (0.07);
  \node[font=\footnotesize, anchor=south east] at (0.95,1.05) {P};
  \node[font=\footnotesize, anchor=south west] at (5.05,4.05) {Q};
  \node[font=\footnotesize, anchor=north west] at (5.05,0.95) {R};
  \node[font=\footnotesize, text=bkBlue!80!black, anchor=north] at (3,0.85) {$x_2-x_1$};
  \node[font=\footnotesize, text=bkOrange!80!black, anchor=west] at (5.15,2.5) {$y_2-y_1$};
  \node[font=\small, text=bkRed!80!black, anchor=south east] at (2.9,2.6) {$d$};
\end{tikzpicture}
\end{minipage}\hfill
\begin{minipage}[c]{0.47\linewidth}
点P$(x_1,\ y_1)$ と点Q$(x_2,\ y_2)$ の距離を $d$ とします。Pから横へ、Qからたてへ線をのばし、交わる点をR$(x_2,\ y_1)$ とします。三角形PRQは、Rが直角の直角三角形です。
\begin{itemize}[leftmargin=1.2\zw, itemsep=1pt, topsep=2pt, label={\color{\bkTT}●}]
\item PR（横）の長さは $|x_2-x_1|$
\item RQ（たて）の長さは $|y_2-y_1|$
\end{itemize}
斜辺がPQなので、三平方の定理から $d^2=|x_2-x_1|^2+|y_2-y_1|^2$。
\end{minipage}

\vspace{4pt}
\bkivehead{絶対値の記号は、なくしてよい}
2乗すると、マイナスの数もプラスになります。たとえば $|-3|^2=3^2=9$、$(-3)^2=9$ で同じです。だから、2乗するときは絶対値の記号 $|\ \ |$ をつけなくてかまいません。両辺の $\sqrt{\phantom{2}}$ をとると（長さは正）、
\bkivewaku{$d=\sqrt{(x_2-x_1)^2+(y_2-y_1)^2}$}
ひく順番を逆にしても、2乗すれば同じです（$(x_1-x_2)^2=(x_2-x_1)^2$）。だから、どちらの点からひいてもかまいません。\bkidx{にてんのきょり}{2点の距離}

\bkivehead{\bkhatenaref の答え}
A$(-2,\ 1)$ と C$(3,\ -2)$ の距離を求めます。
\begin{align*}
x\ \text{の差}&\quad 3-(-2)=5\\
y\ \text{の差}&\quad -2-1=-3\\
\mathrm{AC}&=\sqrt{5^2+(-3)^2}=\sqrt{25+9}=\sqrt{34}=5.830\cdots
\end{align*}
AとCの距離は $\sqrt{34}$ で、約5.83です。ものさしではかった5.8くらいと合っています。「Bを通る道」の8より、約2.2短い近道です。

\bkivehead{地図の上の直線距離}
駅を原点にして、東を $x$、北を $y$ の向き（単位はkm）とします。学校が $(1.2,\ 0.5)$ にあるなら、駅から学校までのまっすぐの距離は
\bkiveline{$\sqrt{1.2^2+0.5^2}=\sqrt{1.44+0.25}=\sqrt{1.69}=1.3$\qquad だから \qty{1.3}{km}}
スマートフォンの地図アプリが表示する「直線距離」は、こうして計算できます。実際に歩く道のりは、道が曲がっている分、もっと長くなります。

\bksubsection{7}{原点を中心とする円の式}
\label{bk:4-4:en}原点Oを中心とする、半径5の円を考えます。円は「中心から同じ距離にある点の集まり」です。だから、円の上の点 $(x,\ y)$ は、Oからの距離がちょうど5です。

\noindent\begin{minipage}[c]{0.46\linewidth}\centering
\begin{tikzpicture}[x=4.0mm, y=4.0mm]
  \draw[bkLine, line width=0.3pt, step=1] (-6,-6) grid (6,6);
  \draw[bkGray, line width=0.6pt, -{Stealth[length=4pt]}] (-6.3,0) -- (6.6,0) node[anchor=west, font=\sffamily\footnotesize] {$x$};
  \draw[bkGray, line width=0.6pt, -{Stealth[length=4pt]}] (0,-6.3) -- (0,6.6) node[anchor=south, font=\sffamily\footnotesize] {$y$};
  \draw[bkBlue, line width=1.3pt] (0,0) circle (5);
  \fill[bkYellow!40] (0,0) -- (3,0) -- (3,4) -- cycle;
  \draw[bkInk, line width=0.9pt] (0,0) -- (3,0) -- (3,4) -- cycle;
  \bkiveright{(3,0)}{(0,0)}{(3,4)}
  \foreach \p in {(3,4),(-4,3),(5,0),(0,-5),(-3,-4)} {\fill[bkRed] \p circle (0.2);}
  \node[font=\footnotesize, anchor=north, fill=white, inner sep=0.8pt] at (1.5,-0.25) {3};
  \node[font=\footnotesize, anchor=west, fill=white, inner sep=0.8pt] at (3.2,1.6) {4};
  \node[font=\footnotesize, anchor=south east, fill=white, inner sep=0.8pt] at (1.3,2.2) {5};
  \node[font=\footnotesize, anchor=south west, fill=white, inner sep=0.8pt] at (3.15,4.15) {$(3,\,4)$};
  \node[font=\footnotesize, anchor=south east, fill=white, inner sep=0.8pt] at (-4.15,3.15) {$(-4,\,3)$};
  \fill[bkGreen!70!black] (-2,1) circle (0.2);
  \node[font=\footnotesize, anchor=north, fill=white, inner sep=0.8pt, text=bkGreen!50!black] at (-2,0.6) {$(-2,\,1)$};
  \draw[bkPurple, line width=0.8pt] (4,-4) circle (0.25);
  \node[font=\footnotesize, anchor=north west, fill=white, inner sep=0.8pt, text=bkPurple] at (4.25,-4.25) {$(4,\,-4)$};
\end{tikzpicture}
\end{minipage}\hfill
\begin{minipage}[c]{0.51\linewidth}
点 $(x,\ y)$ とOの距離は、⑥の式で $(x_1,\ y_1)=(0,\ 0)$ として $\sqrt{x^2+y^2}$ です。これが5なので、両辺を2乗して、
\bkiveline{$x^2+y^2=25$}
赤い点で確かめます。
\begin{itemize}[leftmargin=1.2\zw, itemsep=0pt, topsep=1pt, label={\color{\bkTT}●}]
\item $(3,\ 4)$：$9+16=25$
\item $(-4,\ 3)$：$16+9=25$
\item $(5,\ 0)$：$25+0=25$
\item $(0,\ -5)$：$0+25=25$
\end{itemize}
どれも円の上にあります。
\end{minipage}

\vspace{4pt}
円の内側の点は、Oからの距離が5より短いので、$x^2+y^2$ は25より小さくなります。緑の点 $(-2,\ 1)$ は $4+1=5$ で25より小さく、内側です。紫の点 $(4,\ -4)$ は $16+16=32$ で25より大きく、外側です。

半径を $r$ にしても同じなので、\textbf{原点を中心とする半径 $r$ の円の式}は、
\bkivewaku{$x^2+y^2=r^2$}
です。この $r$ は円の半径です。3-15の相関係数 $r$ とは別の意味なので、気をつけましょう。\bkidx{えんのしき}{円の式}

\bkivehead{円の上半分を、$y=\cdots$ の形で書く}
$x^2+y^2=r^2$ を $y$ について解きます。$x^2$ を右へ移すと $y^2=r^2-x^2$。2乗して $r^2-x^2$ になる数は2つあり（3-10）、$y=\sqrt{r^2-x^2}$ と $y=-\sqrt{r^2-x^2}$ です。上半分（$y$ が0以上）は
\bkiveline{$y=\sqrt{r^2-x^2}$}
$r=5$ で表を作ると、点は円の上半分にきれいに並びます。

\noindent\begin{minipage}[c]{0.5\linewidth}\centering
\begin{bktable}{colspec={Q[c,m,wd=8mm] Q[c,m,wd=20mm] Q[c,m,wd=18mm]}, rowsep=0pt, cells={font=\small}}
$x$ & $25-x^2$ & $y$ \\
0 & 25 & 5 \\
1 & 24 & 4.90 \\
2 & 21 & 4.58 \\
3 & 16 & 4 \\
4 & 9 & 3 \\
5 & 0 & 0 \\
\end{bktable}
\end{minipage}\hfill
\begin{minipage}[c]{0.48\linewidth}\centering
\begin{tikzpicture}
\begin{axis}[bkplot, width=66mm, height=40mm, xmin=-5.8, xmax=5.8, ymin=0, ymax=6.4, xtick={-5,-3,3,5}, ytick={2,4}, xlabel={$x$}, ylabel={$y$}]
  \addplot[bkcurve, domain=-5:5, samples=121] {sqrt(25-x^2)};
  \addplot[only marks, mark=*, mark size=1.6pt, bkOrange] coordinates {(0,5) (1,4.899) (2,4.583) (3,4) (4,3) (5,0) (-1,4.899) (-2,4.583) (-3,4) (-4,3) (-5,0)};
\end{axis}
\end{tikzpicture}
\end{minipage}

\vspace{3pt}
$x$ がマイナスのときも、$x^2$ は同じなので、$y$ は同じ値になります（左右対称）。5-16では、この $y=\sqrt{r^2-x^2}$ を使って、円の面積を計算します。

\bkivehead{中心が原点でない円（軽く）}
\noindent\begin{minipage}[c]{0.58\linewidth}
中心が $(a,\ b)$、半径 $r$ の円の上の点 $(x,\ y)$ は、中心との距離が $r$ なので、⑥の式から
\bkiveline{$(x-a)^2+(y-b)^2=r^2$}
3-9で、$y=x^2$ を右へ $p$ ずらすと $x$ が $x-p$ になったのと同じです。円を右へ $a$、上へ $b$ ずらすと、$x$ が $x-a$、$y$ が $y-b$ になります。

たとえば、中心 $(2,\ 1)$、半径3の円は
\bkiveline{$(x-2)^2+(y-1)^2=9$}
です。点 $(5,\ 1)$ を入れると $9+0=9$ で、円の上にあります。
\end{minipage}\hfill
\begin{minipage}[c]{0.38\linewidth}\centering
\begin{tikzpicture}[x=4.1mm, y=4.1mm]
  \draw[bkLine, line width=0.3pt, step=1] (-3,-3) grid (6,5);
  \draw[bkGray, line width=0.6pt, -{Stealth[length=4pt]}] (-3.3,0) -- (6.6,0) node[anchor=west, font=\sffamily\footnotesize] {$x$};
  \draw[bkGray, line width=0.6pt, -{Stealth[length=4pt]}] (0,-3.3) -- (0,5.6) node[anchor=south, font=\sffamily\footnotesize] {$y$};
  \draw[bkGray!70, dashed, line width=0.6pt] (0,0) circle (3);
  \draw[bkBlue, line width=1.3pt] (2,1) circle (3);
  \fill[bkInk] (2,1) circle (0.15);
  \draw[bkInk, line width=0.6pt] (2,1) -- (5,1);
  \fill[bkRed] (5,1) circle (0.18);
  \node[font=\footnotesize, anchor=south, fill=white, inner sep=0.8pt] at (3.5,1.2) {3};
  \node[font=\footnotesize, anchor=south east, fill=white, inner sep=0.8pt] at (1.85,1.15) {$(2,\,1)$};
  \node[font=\footnotesize, anchor=west, fill=white, inner sep=0.8pt] at (5.2,1.5) {$(5,\,1)$};
\end{tikzpicture}
\par\vspace{1pt}{\sffamily\footnotesize 点線の円（原点が中心）を、\\右へ2、上へ1ずらすと青い円}
\end{minipage}

\bksubsection{8}{直方体の対角線}
\label{bk:4-4:rippou}三平方の定理は、立体の中のななめの長さにも使えます。横 \qty{3}{cm}・おく行き \qty{4}{cm}・高さ \qty{12}{cm} の直方体で、いちばん遠い2つの頂点を結ぶ線分（対角線）の長さを求めます。

\noindent\begin{minipage}[c]{0.4\linewidth}\centering
\begin{tikzpicture}[x={(1mm,0mm)}, y={(0.85mm,0.6mm)}, z={(0mm,1mm)}, scale=3.2]
  \coordinate (O) at (0,0,0); \coordinate (P) at (3,0,0); \coordinate (Q) at (3,4,0); \coordinate (S) at (0,4,0);
  \coordinate (Ot) at (0,0,12); \coordinate (Pt) at (3,0,12); \coordinate (Qt) at (3,4,12); \coordinate (St) at (0,4,12);
  \fill[bkBlue!10] (O) -- (P) -- (Q) -- (S) -- cycle;
  \draw[bkGray, line width=0.5pt, dashed] (O) -- (S) -- (Q) (S) -- (St);
  \fill[bkYellow!30] (O) -- (Q) -- (Qt) -- cycle;
  \draw[bkInk, line width=0.8pt] (O) -- (P) -- (Q) -- (Qt) -- (Pt) -- (P) (O) -- (Ot) -- (Pt) (Ot) -- (St) -- (Qt);
  \draw[bkBlue, line width=1.3pt] (O) -- (Q);
  \draw[bkRed, line width=1.4pt] (O) -- (Qt);
  \node[font=\footnotesize, anchor=north] at (1.5,0,-0.4) {3};
  \node[font=\footnotesize, anchor=north west] at (3.1,2,0) {4};
  \node[font=\footnotesize, anchor=west] at (3.2,4,6) {12};
  \node[font=\footnotesize, text=bkBlue!80!black, anchor=north west, inner sep=0.5pt] at (2.0,1.6,0) {5};
  \node[font=\footnotesize, text=bkRed!80!black, anchor=east, inner sep=0.5pt] at (0.9,1.2,6) {13};
\end{tikzpicture}
\end{minipage}\hfill
\begin{minipage}[c]{0.57\linewidth}
\bkivestep{1}{底面の対角線（青）}
底面は、横3・おく行き4の長方形です。その対角線は、直角をはさむ辺が3と4の直角三角形の斜辺なので、
\bkiveline{$\sqrt{3^2+4^2}=\sqrt{25}=5$}
\bkivestep{2}{立体の対角線（赤）}
底面の対角線（5）と、高さの辺（12）は、直角に交わります（高さの辺は、底面にまっすぐ立っているから）。黄色の三角形は直角三角形なので、
\bkiveline{$\sqrt{5^2+12^2}=\sqrt{169}=13$}
\end{minipage}

\vspace{4pt}
対角線の長さは \qty{13}{cm} です。三平方の定理を2回使いました。①で出てきた $5^2$ は $3^2+4^2$ なので、まとめると
\bkivewaku{横 $a$・おく行き $b$・高さ $c$ の直方体の対角線＝$\sqrt{a^2+b^2+c^2}$}
確かめ：$\sqrt{9+16+144}=\sqrt{169}=13$。4-13では、この考えを、空間の中の2点の距離に広げます。



\Needspace{21\baselineskip}
\begin{motto}[3]
\textbf{三平方の定理をくり返して、$\pi$ を求める}\par
2-9で「12角形から先の数は、4-4で学ぶ三平方の定理をくり返し使うと計算できます」と書いた、その計算です。

\noindent\begin{minipage}[c]{0.36\linewidth}\centering
\begin{tikzpicture}[x=0.8mm, y=0.8mm]
  \draw[bkBlue!60, line width=0.8pt] (60:40) arc (60:120:40);
  \coordinate (O) at (0,0); \coordinate (P) at (60:40); \coordinate (Q) at (120:40); \coordinate (M) at (0,34.64); \coordinate (T) at (0,40);
  \draw[bkInk, line width=0.5pt] (O) -- (P) (O) -- (Q);
  \draw[bkInk, line width=1.0pt] (P) -- (Q);
  \draw[bkGray, dashed, line width=0.6pt] (O) -- (M);
  \draw[bkRed, line width=0.8pt] (M) -- (T);
  \draw[bkOrange!80!black, line width=1.2pt] (P) -- (T);
  \fill[bkInk] (O) circle (0.7);
  \node[font=\footnotesize, anchor=north] at (0,-0.8) {中心};
  \node[font=\footnotesize, anchor=north] at (10,33.8) {$\dfrac{s}{2}$};
  \node[font=\footnotesize, anchor=east] at (-0.6,18) {$h$};
  \node[font=\footnotesize, anchor=south, text=bkOrange!80!black] at (12,39) {新しい辺};
  \node[font=\footnotesize, anchor=west] at (12,17) {1};
\end{tikzpicture}
\end{minipage}\hfill
\begin{minipage}[c]{0.6\linewidth}
半径1の円の内側に、1辺 $s$ の正多角形があるとします。角の数を2倍にした正多角形の1辺（オレンジ）を、三平方の定理で求めます。
\begin{enumerate}[label={\bkcircled[bkGray]{\arabic*}}, leftmargin=2.2\zw, itemsep=0pt, topsep=1pt]
\item 中心から辺までの高さ\par $h=\sqrt{1-\Bigl(\dfrac{s}{2}\Bigr)^2}$
\item 辺から円までのすき間（赤）：$1-h$
\item 新しい辺：$\sqrt{\Bigl(\dfrac{s}{2}\Bigr)^2+(1-h)^2}$
\end{enumerate}
\end{minipage}

\vspace{3pt}
正六角形（$s=1$）から始めて、表計算でこれをくり返すと、次のようになります。
\begin{bkcenter}
\begin{bktable}{colspec={Q[l,m,wd=28mm] *{5}{Q[c,m,wd=14mm]}}, rowsep=1pt, cells={font=\small}}
角の数 & 6 & 12 & 24 & 48 & 96 \\
1辺の長さ & 1 & 0.5176 & 0.2611 & 0.1308 & 0.0654 \\
まわり÷直径 & 3 & 3.1058 & 3.1326 & 3.1394 & 3.1410 \\
\end{bktable}
\end{bkcenter}
2-9の表の「$\pi$ より小さい（内側）」の数と同じになり、3.14159…に近づいていきます。アルキメデスは2200年ほど前に、表計算も電卓もなしで、正96角形まで計算しました。
\end{motto}

\begin{taisetsu}
\def\tsl#1{\makebox[6.6\zw][l]{\sffamily\bfseries\small\color{bkDeep}#1}}
\tsl{三平方の定理}直角三角形で、斜辺を $c$ とすると　$a^2+b^2=c^2$\\[3pt]
\tsl{逆}$a^2+b^2=c^2$ なら、$c$ の向かいの角は直角\\[3pt]
\tsl{三角定規}$1:1:\sqrt2$（45°）　$1:2:\sqrt3$（30°・60°）\\[3pt]
\tsl{2点の距離}$d=\sqrt{(x_2-x_1)^2+(y_2-y_1)^2}$\\[3pt]
\tsl{円の式}原点が中心、半径 $r$ の円は　$x^2+y^2=r^2$\\[3pt]
\tsl{直方体}対角線＝$\sqrt{a^2+b^2+c^2}$
\end{taisetsu}

\begin{kotoba}
\small
\bkword{斜辺}{直角三角形で、直角の向かいにある辺。いちばん長い。}
\bkword{三平方の定理}{直角三角形で $a^2+b^2=c^2$（$c$ は斜辺）。ピタゴラスの定理ともいう。}
\bkword{逆}{「AならばB」に対する「BならばA」。三平方の定理の逆は正しい。}
\bkword{対角線}{多角形や立体で、となり合わない2つの頂点を結ぶ線分。}
\bkword{直線距離}{2点をまっすぐ結んだ長さ。座標から $\sqrt{(x_2-x_1)^2+(y_2-y_1)^2}$ で求める。}
\end{kotoba}

\bkkigocards{\bkkigocard{\sffamily in}{インチ}{長さの単位。1インチは \qty{2.54}{cm}。画面の大きさは対角線の長さで表す}{inch}{4-4}}{}

\begin{column}{3800年前の粘土板}
イラクで見つかった約3800年前の粘土板（プリンプトン322）には、$119^2+120^2=169^2$ のように、$a^2+b^2=c^2$ をみたす数の組が15組も記されています。ピタゴラスより1000年以上前から、人々はこのきまりを知っていたようです。証明の方法は、知られているだけで数百通りもあります。
\end{column}

%% ---------------- 例題 ----------------
\bksection{reidai}{例題でたしかめよう}

\begin{reidai}[斜辺を求める]
直角をはさむ2辺が次の長さの直角三角形で、斜辺の長さを求めましょう。\quad（1）\qty{5}{cm} と \qty{12}{cm}\quad（2）\qty{2}{cm} と \qty{3}{cm}\par
\kangaekata 斜辺を $c$ として、$c^2=a^2+b^2$。$c$ は、2乗するとその数になる正の数です。\par
\shiki（1）$c^2=5^2+12^2=25+144=169$\par
\hspace*{4.6\zw}$c=\sqrt{169}=13$\quad（$13^2=169$ だから）\par
\hspace*{2.6\zw}（2）$c^2=2^2+3^2=4+9=13$\par
\hspace*{4.6\zw}$c=\sqrt{13}=3.605\cdots$\quad（電卓で計算）\par
\kotae（1）\qty{13}{cm}\quad（2）$\sqrt{13}$\,cm（約 \qty{3.61}{cm}）
\end{reidai}

\begin{reidai}[直角をはさむ辺を求める]
斜辺と1つの辺が次の長さの直角三角形で、のこりの辺の長さを求めましょう。\par
（1）斜辺 \qty{13}{cm}、1辺 \qty{5}{cm}\quad（2）斜辺 \qty{6}{cm}、1辺 \qty{2}{cm}\par
\kangaekata $a^2+b^2=c^2$ を、わからない辺について解く（2-14）。$a^2=c^2-b^2$。$\sqrt{\phantom{2}}$ の中はできるだけ小さい整数にする（3-10）。\par
\shiki（1）$a^2=13^2-5^2=169-25=144$\par
\hspace*{4.6\zw}$a=\sqrt{144}=12$\par
\hspace*{2.6\zw}（2）$a^2=6^2-2^2=36-4=32$\par
\hspace*{4.6\zw}$a=\sqrt{32}=\sqrt{16\times2}=4\sqrt2=5.656\cdots$\par
\kotae（1）\qty{12}{cm}\quad（2）$4\sqrt2$\,cm（約 \qty{5.66}{cm}）
\end{reidai}

\noindent\begin{minipage}[t]{0.485\linewidth}\vspace{0pt}\raggedright
\begin{jibun}\raggedright
直角をはさむ2辺が次の長さの直角三角形で、斜辺の長さを求めましょう。\par（1）\qty{8}{cm} と \qty{15}{cm}\par（2）\qty{1}{cm} と \qty{3}{cm}
\bkjans{（1）\nobreak $\sqrt{64+225}=\sqrt{289}=\qty{17}{cm}$。（2）\nobreak $\sqrt{1+9}=\sqrt{10}$\,cm（約 \qty{3.16}{cm}）}
\end{jibun}
\end{minipage}\hfill
\begin{minipage}[t]{0.485\linewidth}\vspace{0pt}\raggedright
\begin{jibun}\raggedright
斜辺と1つの辺が次の長さの直角三角形で、のこりの辺の長さを求めましょう。\par
（1）斜辺 \qty{25}{cm}、1辺 \qty{7}{cm}\par（2）斜辺 \qty{4}{cm}、1辺 \qty{2}{cm}
\bkjans{（1）\nobreak $\sqrt{625-49}=\sqrt{576}=\qty{24}{cm}$。（2）\nobreak $\sqrt{16-4}=\sqrt{12}=2\sqrt3$\,cm（約 \qty{3.46}{cm}）}
\end{jibun}
\end{minipage}

\begin{reidai}[三角定規の比を使う]
（1）1辺 \qty{6}{cm} の正方形の対角線の長さを求めましょう。\par
（2）1辺 \qty{8}{cm} の正三角形の高さを求めましょう。\par
\kangaekata（1）対角線で切ると直角二等辺三角形で、$1:1:\sqrt2$。（2）高さで半分にすると、30°・60°・90°の直角三角形で、短い辺：斜辺：のこり＝$1:2:\sqrt3$。\par
\shiki（1）$1:\sqrt2=6:(\text{対角線})$ なので、対角線は $6\times\sqrt2=6\sqrt2$\par
\hspace*{2.6\zw}（2）斜辺8、短い辺は $8\div2=4$\par
\hspace*{4.6\zw}$1:\sqrt3=4:(\text{高さ})$ なので、高さは $4\times\sqrt3=4\sqrt3$\par
\hspace*{4.6\zw}確かめ：$\sqrt{8^2-4^2}=\sqrt{48}=\sqrt{16\times3}=4\sqrt3$\par
\kotae（1）$6\sqrt2$\,cm（約 \qty{8.49}{cm}）\quad（2）$4\sqrt3$\,cm（約 \qty{6.93}{cm}）
\end{reidai}

\begin{reidai}[直角三角形かどうかを見分ける]
3辺が次の長さの三角形は、直角三角形ですか。\quad（1）9、12、15\quad（2）4、5、6\par
\kangaekata いちばん長い辺を $c$ にして、$a^2+b^2$ と $c^2$ をくらべる。等しければ、三平方の定理の逆から直角三角形（⑤）。\par
\shiki（1）$9^2+12^2=81+144=225$、$15^2=225$。等しい\par
\hspace*{2.6\zw}（2）$4^2+5^2=16+25=41$、$6^2=36$。等しくない\par
\kotae（1）直角三角形（15の辺の向かいが直角）\quad（2）直角三角形ではない
\end{reidai}

\noindent\begin{minipage}[t]{0.485\linewidth}\vspace{0pt}\raggedright
\begin{jibun}\raggedright
（1）1辺 \qty{10}{cm} の正方形の対角線の長さを求めましょう。\par
（2）1辺 \qty{6}{cm} の正三角形の高さを求めましょう。
\bkjans{（1）\nobreak $10\sqrt2$\,cm（約 \qty{14.1}{cm}）。（2）\nobreak $3\sqrt3$\,cm（約 \qty{5.20}{cm}）}
\end{jibun}
\end{minipage}\hfill
\begin{minipage}[t]{0.485\linewidth}\vspace{0pt}\raggedright
\begin{jibun}\raggedright
3辺が次の長さの三角形は、直角三角形ですか。\par（1）7、24、25\par（2）5、6、8
\bkjans{（1）\nobreak $49+576=625=25^2$ なので、直角三角形。（2）\nobreak $25+36=61$、$8^2=64$ で等しくないので、直角三角形ではない}
\end{jibun}
\end{minipage}

\begin{reidai}[座標の上の2点の距離]
次の2点の距離を求めましょう。\quad（1）A$(1,\ 2)$、B$(4,\ 6)$\quad（2）P$(-1,\ 3)$、Q$(3,\ -5)$\par
\kangaekata $d=\sqrt{(x_2-x_1)^2+(y_2-y_1)^2}$（⑥）。マイナスの数は、かっこをつけて入れる。\par
\shiki（1）$x$ の差 $4-1=3$、$y$ の差 $6-2=4$\par
\hspace*{4.6\zw}$\sqrt{3^2+4^2}=\sqrt{25}=5$\par
\hspace*{2.6\zw}（2）$x$ の差 $3-(-1)=4$、$y$ の差 $-5-3=-8$\par
\hspace*{4.6\zw}$\sqrt{4^2+(-8)^2}=\sqrt{16+64}=\sqrt{80}=\sqrt{16\times5}=4\sqrt5$\par
\kotae（1）5\quad（2）$4\sqrt5$（約8.94）
\end{reidai}

\begin{reidai}[原点を中心とする円の式]
（1）原点を中心とし、点 $(6,\ 8)$ を通る円の式を求めましょう。\par
（2）点 $(2,\ 4)$ は、円 $x^2+y^2=25$ の上・内側・外側のどれにありますか。\par
\kangaekata（1）半径は、原点と $(6,\ 8)$ の距離。$x^2+y^2=r^2$ の $r^2$ を求めればよい。（2）$x^2+y^2$ を計算して、$r^2=25$ とくらべる（⑦）。\par
\noindent\begin{minipage}[c]{0.6\linewidth}
\shiki（1）$r^2=6^2+8^2=36+64=100$\par
\hspace*{2.6\zw}（2）$2^2+4^2=4+16=20$\par
\hspace*{4.6\zw}20は25より小さい\par
\kotae（1）$x^2+y^2=100$（半径10）\par
\hspace*{2.6\zw}（2）内側
\end{minipage}\hfill
\begin{minipage}[c]{0.38\linewidth}\centering
\begin{tikzpicture}[x=1.4mm, y=1.4mm]
  \draw[bkGray, line width=0.5pt, -{Stealth[length=3pt]}] (-11.5,0) -- (12,0) node[anchor=west, font=\sffamily\footnotesize] {$x$};
  \draw[bkGray, line width=0.5pt, -{Stealth[length=3pt]}] (0,-11.5) -- (0,12) node[anchor=south, font=\sffamily\footnotesize] {$y$};
  \draw[bkBlue, line width=1.1pt] (0,0) circle (10);
  \fill[bkYellow!40] (0,0) -- (6,0) -- (6,8) -- cycle;
  \draw[bkInk, line width=0.7pt] (0,0) -- (6,0) -- (6,8) -- cycle;
  \fill[bkRed] (6,8) circle (0.45);
  \node[font=\footnotesize, anchor=north] at (3,-0.3) {6};
  \node[font=\footnotesize, anchor=west] at (6.3,3.5) {8};
  \node[font=\footnotesize, anchor=south east] at (2.8,4.2) {10};
  \node[font=\footnotesize, anchor=south west] at (6.3,8.3) {$(6,\,8)$};
\end{tikzpicture}
\end{minipage}
\end{reidai}

\noindent\begin{minipage}[t]{0.485\linewidth}\vspace{0pt}\raggedright
\begin{jibun}\raggedright
次の2点の距離を求めましょう。\par（1）$(2,\ 1)$、$(8,\ 9)$\par（2）$(-4,\ 1)$、$(2,\ -3)$
\bkjans{（1）\nobreak $\sqrt{36+64}=\sqrt{100}=10$。（2）\nobreak $\sqrt{6^2+(-4)^2}=\sqrt{52}=2\sqrt{13}$（約7.21）}
\end{jibun}
\end{minipage}\hfill
\begin{minipage}[t]{0.485\linewidth}\vspace{0pt}\raggedright
\begin{jibun}\raggedright
（1）原点を中心とし、点 $(5,\ 12)$ を通る円の式を求めましょう。\par
（2）点 $(-8,\ 6)$ は、円 $x^2+y^2=100$ の上・内側・外側のどれにありますか。
\bkjans{（1）\nobreak $r^2=25+144=169$ なので $x^2+y^2=169$。（2）\nobreak $64+36=100$ なので、円の上}
\end{jibun}
\end{minipage}

\begin{reidai}[直方体の対角線]
横 \qty{3}{cm}、おく行き \qty{4}{cm}、高さ \qty{12}{cm} の直方体の対角線の長さを求めましょう。\par
\kangaekata 底面の対角線を求めてから、高さと組にして、もう一度三平方の定理（⑧）。まとめて $\sqrt{a^2+b^2+c^2}$ でもよい。\par
\shiki 底面の対角線：$\sqrt{3^2+4^2}=\sqrt{25}=5$\par
\hspace*{2.6\zw}立体の対角線：$\sqrt{5^2+12^2}=\sqrt{25+144}=\sqrt{169}=13$\par
\kotae \qty{13}{cm}
\end{reidai}

\begin{jibun}
横 \qty{2}{cm}、おく行き \qty{3}{cm}、高さ \qty{6}{cm} の直方体の対角線の長さを求めましょう。
\bkjans{$\sqrt{4+9+36}=\sqrt{49}=\qty{7}{cm}$}
\end{jibun}

\begin{reidai}[かべに立てかけたはしご]
長さ \qty{5}{m} のはしごを、かべに立てかけます。はしごの下のはしを、かべから \qty{1.4}{m} はなすと、上のはしは地面から何mの高さに届きますか。\par
\kangaekata かべと地面は直角。はしごが斜辺になる。図をかいて、どこが斜辺かを確かめる。\par
\noindent\begin{minipage}[c]{0.66\linewidth}
\shiki 高さを $h$ とすると、$h^2+1.4^2=5^2$\par
\hspace*{2.6\zw}$h^2=25-1.96=23.04$\par
\hspace*{2.6\zw}$h=\sqrt{23.04}=4.8$\quad（$4.8^2=23.04$）\par
\kotae \qty{4.8}{m}
\end{minipage}\hfill
\begin{minipage}[c]{0.3\linewidth}\centering
\begin{tikzpicture}[x=1mm, y=1mm]
  \fill[bkGray!30] (-2,-2) rectangle (22,0);
  \fill[bkGray!30] (16,0) rectangle (19,26);
  \draw[bkInk, line width=0.6pt] (-2,0) -- (16,0) -- (16,26);
  \draw[bkOrange!80!black, line width=1.6pt] (9.28,0) -- (16,24);
  \bkiveright{(16,0)}{(9.28,0)}{(16,24)}
  \node[font=\footnotesize, anchor=south east] at (11.8,12) {\qty{5}{m}};
  \node[font=\footnotesize, anchor=north] at (12.6,-2.4) {\qty{1.4}{m}};
  \node[font=\footnotesize, anchor=west] at (19.5,12) {$h$};
\end{tikzpicture}
\end{minipage}
\end{reidai}

\begin{jibun}
長さ \qty{6.5}{m} のはしごの下のはしを、かべから \qty{2.5}{m} はなして立てかけます。上のはしは、地面から何mの高さに届きますか。
\bkjans{$h^2=6.5^2-2.5^2=42.25-6.25=36$、$h=\qty{6}{m}$}
\end{jibun}

\begin{reidai}[画面の大きさ（インチ）]
タブレットの画面は、たて \qty{16}{cm}、横 \qty{12}{cm} の長方形です。画面の大きさ（対角線の長さ）は何インチですか。四捨五入して小数第1位まで求めましょう。\qty{1}{in}＝\qty{2.54}{cm} です。\par
\kangaekata 対角線は、たてと横を直角をはさむ辺とする直角三角形の斜辺。cmで求めてから、単位つきの1（\qty{2.54}{cm/in}）でわってインチに換算する（2-10）。\par
\shiki 対角線：$\sqrt{16^2+12^2}=\sqrt{256+144}=\sqrt{400}=20$\quad（cm）\par
\hspace*{2.6\zw}$\qty{20}{cm}\div\qty{2.54}{cm/in}=\qty{7.87}{in}$\par
\kotae 約7.9インチ
\end{reidai}

\begin{jibun}
たて \qty{24}{cm}、横 \qty{18}{cm} の画面の大きさは何インチですか。四捨五入して小数第1位まで求めましょう。
\bkjans{$\sqrt{576+324}=\sqrt{900}=\qty{30}{cm}$。$30\div2.54=11.81\cdots$ で、約11.8インチ}
\end{jibun}

%% ---------------- まちがえやすいところ ----------------
\bksection{tsumazuki}{まちがえやすいところをチェック}

\begin{tsumazuki}{マイナスの差を2乗するとき}
\bkng{$(-3)^2=-9$}
\bkok{$(-3)^2=(-3)\times(-3)=9$}
マイナス×マイナスはプラスです（2-2）。$-3^2$ と書くと「$3^2$ にマイナスをつけた $-9$」になってしまうので、マイナスの数は必ずかっこをつけて2乗します。
\end{tsumazuki}

\begin{tsumazuki}{斜辺でない辺を $c$ にする}
直角三角形で、直角をはさむ辺の1つが3、斜辺が5のとき、のこりの辺は？\par
\bkng{$3^2+5^2=34$ なので、$\sqrt{34}$}
\bkok{$5^2-3^2=16$ なので、$\sqrt{16}=4$}
$c$ はいつも斜辺（直角の向かいの、いちばん長い辺）です。図をかいて、直角の向かいの辺に印をつけてから式を立てましょう。
\end{tsumazuki}

\begin{tsumazuki}{$\sqrt{\phantom{2}}$ の中のたし算を、ばらばらにする}
\bkng{$\sqrt{3^2+4^2}=3+4=7$}
\bkok{$\sqrt{3^2+4^2}=\sqrt{25}=5$}
$\sqrt{\phantom{2}}$ の中にたし算があるときは、先に中を計算します。7になるのは、Bを通る遠回りの道のりです。まっすぐの距離は、それより短くなります。
\end{tsumazuki}

\begin{tsumazuki}{直角かどうかわからない三角形に使う}
2辺が4と6の三角形の、3つめの辺は？\par
\bkng{$\sqrt{4^2+6^2}=\sqrt{52}$}
\bkok{直角があるとわからなければ、三平方の定理は使えない}
①で見たように、間の角が変わると、3つめの辺も変わります。三平方の定理は、直角三角形のときだけのきまりです。
\end{tsumazuki}

\begin{tsumazuki}{$1:2:\sqrt3$ の斜辺をまちがえる}
30°・60°・90°の三角定規で、いちばん短い辺が \qty{5}{cm} のとき、斜辺は？\par
\bkng{$5\times\sqrt3=5\sqrt3$（約 \qty{8.66}{cm}）}
\bkok{$5\times2=\qty{10}{cm}$}
$1:2:\sqrt3$ は「短い辺：斜辺：のこりの辺」の順です。斜辺はいちばん長い辺なので、2にあたります。$\sqrt3=1.73\cdots$ は2より小さいので、斜辺にはなりません。④の図で、どの辺が2かを確かめましょう。
\end{tsumazuki}

%% ---------------- 物理・工学 ----------------
\begin{butsuri}
\par\addvspace{3pt}\textbf{ロボットの手は、届く？}\par\noindent
\begin{minipage}[c]{0.6\linewidth}
2本のうでの長さが \qty{0.35}{m} と \qty{0.25}{m} のロボットは、肩から最大で
\bkiveline{$0.35+0.25=\qty{0.6}{m}$}
まで手をのばせます。肩を原点にして、物が $(0.3,\ 0.4)$（m）にあるとき、肩からの距離は
\bkiveline{$\sqrt{0.3^2+0.4^2}=\sqrt{0.25}=\qty{0.5}{m}$}
です。0.6 m より短いので、手が届きます。届くとわかったら、関節を何度曲げればよいかを、4-7の余弦定理で計算します（4-0の逆運動学）。
\end{minipage}\hfill
\begin{minipage}[c]{0.36\linewidth}\centering
\begin{tikzpicture}[x=0.85mm, y=0.85mm]
  \fill[bkBlue!8] (0,0) -- (36,0) arc (0:90:36) -- cycle;
  \draw[bkBlue!60, dashed, line width=0.6pt] (36,0) arc (0:90:36);
  \draw[bkGray, line width=0.5pt] (0,0) -- (40,0) (0,0) -- (0,38);
  \draw[bkInk!70, line width=2.4pt, line cap=round] (0,0) -- (3.36,20.73) -- (18,24);
  \fill[bkInk] (0,0) circle (1.4) (3.36,20.73) circle (1);
  \fill[bkOrange] (18,24) circle (1.4);
  \draw[bkRed, line width=0.8pt, dashed] (0,0) -- (18,24);
  \node[font=\footnotesize, anchor=west] at (19.5,24.5) {物};
  \node[font=\footnotesize, anchor=north] at (2,-1) {肩};
  \node[font=\footnotesize, text=bkBlue!80!black, anchor=west] at (26,27) {届く\,範囲};
\end{tikzpicture}
\end{minipage}
\end{butsuri}

\begin{butsuri}
\par\addvspace{3pt}\textbf{画面の「インチ」は対角線}\par\noindent
\begin{minipage}[c]{0.64\linewidth}
テレビやスマートフォンの画面の大きさは、対角線の長さをインチで表します。「50型」のテレビは、対角線が50インチ（約 \qty{127}{cm}）です。いまのテレビの画面は、横：たて＝$16:9$ です。横16・たて9の長方形の対角線は $\sqrt{16^2+9^2}=\sqrt{337}=18.36\cdots$ なので、50型の横の長さは $50\times\dfrac{16}{18.36}=43.6$ インチ、約 \qty{111}{cm} です。テレビ台を選ぶときは、対角線ではなく、この横の長さを見ます。
\end{minipage}\hfill
\begin{minipage}[c]{0.32\linewidth}\centering
\begin{tikzpicture}[x=0.85mm, y=0.85mm]
  \fill[bkInk!80] (-1.5,-1.5) rectangle (41.5,24);
  \fill[bkBlue!25] (0,0) rectangle (40,22.5);
  \draw[bkRed, line width=1.2pt] (0,0) -- (40,22.5);
  \fill[bkInk!70] (16,-1.5) rectangle (24,-5);
  \fill[bkInk!70] (10,-5) rectangle (30,-6.5);
  \node[font=\footnotesize, text=bkRed!80!black, anchor=south] at (20,25.5) {対角線が50インチ};
\end{tikzpicture}
\end{minipage}

\end{butsuri}

\begin{butsuri}
\par\addvspace{3pt}\textbf{地震の震源までの距離}\par\noindent
\begin{minipage}[c]{0.64\linewidth}
地震が起きた地下の場所を\textbf{震源}、その真上の地面の場所を\textbf{震央}といいます。深さ \qty{10}{km} の震源で地震が起き、震央から \qty{24}{km} はなれた町で、ゆれを観測したとします。震源と町の距離は、
\bkiveline{$\sqrt{10^2+24^2}=\sqrt{676}=\qty{26}{km}$}
最初の小さなゆれ（P波）は1秒に約 \qty{7}{km} 進むので、町には約3.7秒で届きます。緊急地震速報は、この計算を一瞬で行います。
\end{minipage}\hfill
\begin{minipage}[c]{0.33\linewidth}\centering
\begin{tikzpicture}[x=0.75mm, y=0.75mm]
  \fill[bkOrange!14] (-3,-14) rectangle (44,0);
  \draw[bkGreen!60!black, line width=1.2pt] (-3,0) -- (44,0);
  \draw[bkGray, dashed, line width=0.6pt] (0,0) -- (0,-10);
  \draw[bkRed, line width=1.2pt] (0,-10) -- (36,0);
  \fill[bkRed] (0,-10) circle (1);
  \fill[bkInk] (0,0) circle (0.8) (36,0) circle (0.8);
  \node[font=\footnotesize, anchor=south] at (0,1) {震央};
  \node[font=\footnotesize, anchor=south] at (36,1) {町};
  \node[font=\footnotesize, anchor=north] at (3,-11.2) {震源};
  \node[font=\footnotesize, anchor=east] at (-0.8,-5) {10};
  \node[font=\footnotesize, anchor=south, fill=white, inner sep=0.8pt] at (18,0.8) {24};
  \node[font=\footnotesize, text=bkRed!80!black, anchor=north west] at (19,-5.5) {26};
  \bkiveright{(0,0)}{(36,0)}{(0,-10)}
\end{tikzpicture}
\end{minipage}

\end{butsuri}




%% ---------------- 章末問題 ----------------
\bkendproblems
\bkproblemgroup{基本}{3}

\bkproblem 直角をはさむ2辺が次の長さの直角三角形で、斜辺の長さを求めましょう（\bkreidairef{1}）。\quad（1）\qty{9}{cm} と \qty{12}{cm}\quad（2）\qty{1}{cm} と \qty{2}{cm}
\bkpans{（1）\nobreak $\sqrt{81+144}=\sqrt{225}=\qty{15}{cm}$。（2）\nobreak $\sqrt5$\,cm（約 \qty{2.24}{cm}）}

\bkproblem 斜辺と1つの辺が次の長さの直角三角形で、のこりの辺の長さを求めましょう（\bkreidairef{2}）。\quad（1）斜辺 \qty{17}{cm}、1辺 \qty{8}{cm}\quad（2）斜辺 \qty{3}{cm}、1辺 \qty{1}{cm}
\bkpans{（1）\nobreak $\sqrt{289-64}=\sqrt{225}=\qty{15}{cm}$。（2）\nobreak $\sqrt{9-1}=\sqrt8=2\sqrt2$\,cm（約 \qty{2.83}{cm}）}

\bkproblem （1）1辺 \qty{4}{cm} の正方形の対角線と、（2）1辺 \qty{10}{cm} の正三角形の高さを求めましょう（\bkreidairef{3}）。
\bkpans{（1）\nobreak $4\sqrt2$\,cm（約 \qty{5.66}{cm}）。（2）\nobreak $5\sqrt3$\,cm（約 \qty{8.66}{cm}）}

\bkproblem 3辺が次の長さの三角形は、直角三角形ですか（\bkreidairef{4}）。\quad（1）8、15、17\quad（2）6、7、9
\bkpans{（1）\nobreak $64+225=289=17^2$ なので、直角三角形。（2）\nobreak $36+49=85$、$9^2=81$ で等しくないので、直角三角形ではない}

\bkproblem 次の2点の距離を求めましょう（\bkreidairef{5}）。\quad（1）$(0,\ 0)$、$(5,\ 12)$\quad（2）$(-3,\ 4)$、$(3,\ -4)$
\bkpans{（1）\nobreak $\sqrt{25+144}=13$。（2）\nobreak $\sqrt{6^2+(-8)^2}=\sqrt{100}=10$}

\bkproblem （1）原点を中心とし、点 $(-5,\ 12)$ を通る円の式を求めましょう。（2）点 $(10,\ -5)$ は、（1）の円の上・内側・外側のどれにありますか（\bkreidairef{6}）。
\bkpans{（1）\nobreak $r^2=25+144=169$ なので $x^2+y^2=169$。（2）\nobreak $100+25=125$ は169より小さいので、内側}

\bkproblem 横 \qty{4}{cm}、おく行き \qty{4}{cm}、高さ \qty{7}{cm} の直方体の対角線の長さを求めましょう（\bkreidairef{7}）。
\bkpans{$\sqrt{16+16+49}=\sqrt{81}=\qty{9}{cm}$}

\bkproblemgroup{チャレンジ}{3}

\bkproblem 長さ \qty{2.6}{m} のはしごの下のはしを、かべから \qty{1}{m} はなして立てかけます。上のはしは、地面から何mの高さに届きますか（\bkreidairef{8}）。
\bkpans{$h^2=2.6^2-1^2=6.76-1=5.76$、$h=\qty{2.4}{m}$}

\bkproblem 横 \qty{32}{cm}、たて \qty{18}{cm} のパソコンの画面の大きさは何インチですか。四捨五入して小数第1位まで求めましょう（\bkreidairef{9}）。
\bkpans{$\sqrt{1024+324}=\sqrt{1348}=36.71\cdots$\,cm。$36.71\div2.54=14.45\cdots$ で、約14.5インチ}

\bkproblem 深さ \qty{16}{km} の震源で地震が起きました。震央から \qty{30}{km} はなれた町と震源の距離を求めましょう（\bkreidairef{1}・\bkreidairef{8}）。
\bkpans{$\sqrt{16^2+30^2}=\sqrt{256+900}=\sqrt{1156}=\qty{34}{km}$}

\begin{furikaeri}
\begin{checklist}
\item 方眼の面積から $a^2+b^2=c^2$ を見つけ、2通りに説明できる。\bkfurref{bk:4-4:shoumei}
\item 辺の長さ・2点の距離・直方体の対角線を求められる。\bkfurref{bk:4-4:kyori}
\item 直角三角形かを見分け、円の式のわけがわかる。\bkfurref{bk:4-4:en}
\end{checklist}
\end{furikaeri}

\bkfukidashi{sora}{たてと横がわかれば、ななめもわかる。次は面積と体積だよ。}
